\chapter{Forensic and Medical Assessment}
\label{chap:forensic-and-medical-assessment}

\classification{Evidence status: forensic and medical evidence review. Confidence baseline: High for page-mapped observations and specimen receipt, Moderate where OCR distortion or medical interpretation limits precision, Low where testimony concerns conditional possibilities or later unresolved forensic work. This chapter does not evaluate competing manner-of-death or intelligence hypotheses, motive, or final cause of death.}

\section{Medical Overview}

The forensic-medical layer contains 38 medical observations, 14 toxicology entries, 7 time-of-death entries, 10 forensic identifier entries, 10 medical uncertainty entries, and 26 priority medical transcription records. The purpose of this chapter is to state what the medical record documents, what it interprets, and what remains unresolved.

The core medical witnesses are Dwyer, Cowan, Bennett, Cleland, Hicks, Lawson, and Durham. Their evidence is not all the same type. Dwyer gave postmortem evidence. Cowan gave analytical toxicology evidence. Bennett gave an early examination and time estimate. Cleland and Hicks gave expert review and interpretation. Lawson and Durham provided postmortem identification-material evidence.

The controlling medical position is unresolved cause. Dwyer recorded heart failure as the immediate medical event but could not state what caused it. Cowan found no signs of common poison in submitted specimens. Cleland and Hicks considered non-natural death and possible poison mechanisms, but the acquired record does not detect a specific poison.

\begin{table}[htbp]
\centering
\caption{Forensic-medical register summary}
\label{tab:forensic-medical-register-summary}
\begin{tabularx}{\textwidth}{>{\raggedright\arraybackslash}p{0.36\textwidth}r>{\raggedright\arraybackslash}X}
\toprule
Register & Count & Use in this chapter \\
\midrule
Medical observation register & 38 & External, internal, toxicological, timing, identifier, and uncertainty observations. \\
Toxicology register & 14 & Substances or classes considered, tested, not detected, discounted, or unresolved. \\
Time-of-death register & 7 & Medical timing estimates and limitations. \\
Forensic identifiers register & 10 & Dental, fingerprint, photographic, death-mask, physical, and biological identifiers. \\
Medical uncertainty register & 10 & Unresolved medical questions and future source needs. \\
Medical transcription index & 26 & Priority page-level medical transcription controls. \\
\bottomrule
\end{tabularx}
\end{table}

\section{External Examination}

Dwyer's external examination described a tallish male, estimated about 45 years, with greying hair and good physical condition. Fingernails and feet appeared cared for, and nails were carefully trimmed. These are physical observations, not evidence of occupation, social status, or lifestyle by themselves.

External trauma evidence is limited. Dwyer recorded no external markings of note and saw no evidence of hypodermic needle use. He also noted two recent right-hand abrasions that did not appear significant. The medical register cautions that this does not equal a complete trauma or injection exclusion without the full autopsy file and body map.

Rigor and lividity are documented. Dwyer described intense postmortem rigidity and deep lividity behind or above the ears and neck. Cleland later discussed the lividity pattern as perhaps surprising but explainable by head support. These observations are relevant to body-position and time-of-death assessment, but they do not by themselves prove where death occurred.

Other external observations include small or uneven pupils of approximately the same size, dried saliva near the right of the mouth, sand in the hair but not in the nostrils or mouth, leg sunburn distribution, and missing teeth. Each item is recorded as an observation. Interpretation remains constrained by the absence of the full autopsy report and original body photographs.

\section{Internal Examination}

Dwyer recorded normal scalp, skull, and brain except for visible small vessels and congestion. He recorded congestion in the pharynx, superficial whitening and a patch of ulceration in the gullet, and a deeply congested stomach with superficial redness and small submucosal haemorrhages.

The stomach evidence is important but limited. Blood was mixed with food in the stomach, and Dwyer considered that the blood was present before postmortem. He estimated food had been in the stomach for up to three or four hours before death, with an anxiety caveat. This does not provide a precise final-meal time.

The organ findings were dominated by congestion. Kidneys, liver, lungs, spleen, and upper vessels were described with congestion or engorgement. The heart was normal in size and otherwise normal, with firm or tough muscle and a contracted appearance. The spleen was described as strikingly large and firm, about three times normal size.

Microscopy reportedly showed destruction of centres of liver lobules, but the register treats this as partially documented because OCR damage and later interpretation limit precision. Cleland later did not find the liver section explanatory. The acquired record therefore does not convert liver findings into a final cause.

\begingroup
\scriptsize
\setlength{\tabcolsep}{2pt}
\renewcommand{\arraystretch}{1.12}
\begin{longtable}{>{\raggedright\arraybackslash}p{0.12\textwidth}>{\raggedright\arraybackslash}p{0.18\textwidth}>{\raggedright\arraybackslash}p{0.34\textwidth}>{\raggedright\arraybackslash}p{0.16\textwidth}>{\raggedright\arraybackslash}p{0.12\textwidth}}
\caption{Selected medical observations}\label{tab:chapter9-medical-observations}\\
\toprule
ID & Category & Observation & Source & Confidence \\
\midrule
\endfirsthead
\toprule
ID & Category & Observation & Source & Confidence \\
\midrule
\endhead
MED-0001 & External condition & Tallish male, about 45 years, greying hair, good physical condition. & Dwyer p.35 & High \\
MED-0002 & External condition & Fingernails and feet appeared cared for. & Dwyer p.35 & High \\
MED-0003 & Injuries/trauma & No external markings of note recorded. & Dwyer p.35 & High \\
MED-0004 & Rigor/lividity & Intense rigidity and deep lividity behind or above ears and neck. & Dwyer p.35 & High \\
MED-0005 & Dental & Several teeth missing; missing-teeth chart prepared. & Dwyer p.35; DOC-0017 & High \\
MED-0008 & External condition & Sand in hair but not nostrils or mouth. & Dwyer p.37 & High \\
MED-0011 & Stomach & Deep congestion, redness, and small submucosal haemorrhages. & Dwyer p.37 & High \\
MED-0012 & Stomach & Blood mixed with food in stomach. & Dwyer pp.37--43 & High \\
MED-0013 & Organ congestion & Kidneys, liver, lungs, spleen, and upper vessels congested or engorged. & Dwyer pp.37--39 & High \\
MED-0014 & Heart & Heart normal size and otherwise normal; firm muscle; contracted appearance. & Dwyer p.37 & High \\
MED-0018 & Uncertainty & Immediate heart failure recorded, cause of heart failure not stated. & Dwyer p.39 & High \\
\bottomrule
\end{longtable}
\endgroup

\section{Toxicology}

Cowan received stomach contents, liver and muscle, urine, and blood on 2 December 1948. He analyzed submitted specimens and found no signs of any common poison. The chapter uses the controlled wording no common poison detected. It does not extend that result into a broader absence claim.

The toxicology register records common poisons, cyanides, alkaloids, barbiturates, soluble hypnotic, carbolic acid, morphine, insulin, diphtheria toxin, botulism toxin, nicotine, aconite or aconitine class, curare or tubarine, and an unspecified cardiac glycoside or drug group. Several were considered by medical witnesses. None is recorded as detected in the acquired primary text.

The strongest toxicology fact is negative common-poison testing in submitted specimens. The main limitation is that the exact analytical methods, full panel, sensitivity, sample condition, and laboratory notes are not acquired. Cowan thought that if death had been caused by a common poison, his examination would reveal it; if poison caused death, he considered a very rare poison possible.

Dwyer raised a barbiturate or soluble-hypnotic possibility. He later accepted Cowan's negative findings while preserving possibility language. Cleland and Hicks discussed uncommon or difficult-to-detect possibilities. Hicks's specific drug-group or exhibit terminology remains obscured in OCR and must not be named from the current reproduction.

\begingroup
\scriptsize
\setlength{\tabcolsep}{2pt}
\renewcommand{\arraystretch}{1.12}
\begin{longtable}{>{\raggedright\arraybackslash}p{0.12\textwidth}>{\raggedright\arraybackslash}p{0.18\textwidth}>{\raggedright\arraybackslash}p{0.22\textwidth}>{\raggedright\arraybackslash}p{0.16\textwidth}>{\raggedright\arraybackslash}p{0.20\textwidth}}
\caption{Toxicology summary}\label{tab:chapter9-toxicology-summary}\\
\toprule
ID & Substance or class & Evidence position & Detected & Current status \\
\midrule
\endfirsthead
\toprule
ID & Substance or class & Evidence position & Detected & Current status \\
\midrule
\endhead
TOX-0001 & Common poisons & Cowan analyzed submitted specimens and found no signs. & No & Not detected \\
TOX-0002 & Cyanides & Listed or discussed; rapidity and absence of smell/bottle noted. & No & Not detected \\
TOX-0003 & Alkaloids & Listed or discussed as broad class. & No & Not detected \\
TOX-0004 & Barbiturates & Suggested and discussed; Cowan found none; timing/pathology problematic. & No & Considered but unproven \\
TOX-0005 & Soluble hypnotic & Suggested by Dwyer without detection. & No & Suggested by witness \\
TOX-0007 & Morphine & Discussed by Hicks and considered unlikely. & No & Ruled unlikely \\
TOX-0008 & Insulin & Considered; no specific test documented; medically discounted in testimony. & No & Ruled unlikely \\
TOX-0013 & Curare/tubarine & Considered; injection/detection limits unresolved. & No & Considered but unproven \\
TOX-0014 & Unspecified cardiac drug group & Hicks suspected difficult-to-detect group; OCR obscures names. & No & Suggested by witness \\
\bottomrule
\end{longtable}
\endgroup

\section{Time of Death}

The master medical timing record begins with Bennett. Bennett examined the body at 9:40 a.m. on 1 December 1948 and estimated death could have occurred up to eight hours earlier, not more than eight hours. The register places the estimate around about 2:00 a.m., with moderate confidence.

The limitations are significant. Bennett's estimate was based on a cursory look, rigor mortis, and cyanosis, and he had not recorded the extent of rigor at the time. Dwyer's stomach-content timing is a physiological estimate with an anxiety caveat, not an observed time of death.

Hicks and Cleland offered conditional interpretive timing comments. Hicks discussed a seven-hour interval from the 7:00 p.m. sighting to a 2:00 a.m. death estimate only as conditional reasoning dependent on an unestablished toxic agent and unestablished identity of the man seen at 7:00 p.m. Cleland gave a more uncertain estimate in one passage that conflicts with Bennett and is not adopted as the master chronology event.

\begin{table}[htbp]
\centering
\caption{Time-of-death assessment}
\label{tab:chapter9-time-of-death}
\begin{tabularx}{\textwidth}{>{\raggedright\arraybackslash}p{0.14\textwidth}>{\raggedright\arraybackslash}p{0.22\textwidth}>{\raggedright\arraybackslash}p{0.22\textwidth}>{\raggedright\arraybackslash}X}
\toprule
ID & Source & Timing statement & Limitation \\
\midrule
TOD-0001 & Coroner remarks / Bennett & Death around about 2:00 a.m. & Secondary within inquest remarks; exact basis in Bennett testimony. \\
TOD-0002 & Bennett & Up to eight hours before 9:40 a.m.; not more than eight hours. & Cursory look; rigor extent not recorded. \\
TOD-0003 & Dwyer & Food in stomach up to three or four hours before death. & Digestion variable; anxiety caveat. \\
TOD-0005 & Hicks & Conditional seven-hour pharmacological reasoning. & Depends on unestablished toxic-agent and sighting-identity assumptions. \\
TOD-0007 & Cleland & At or before midnight in one passage. & Very low confidence; conflicts with Bennett. \\
\bottomrule
\end{tabularx}
\end{table}

\section{Identification Material}

The forensic identifiers register records dental charting, dental characteristics, fingerprints, photographs, physical measurements, scars or marks, death mask or bust, biological samples, and the death certificate. These materials support identification attempts and future comparison, not cause-of-death conclusions.

Dwyer's dental chart is a strong existence record. Several teeth were missing, and the chart reproduction is preserved as DOC-0017. Exact dental notation remains unresolved because the chart image has no OCR and requires specialist or manual transcription.

Durham fingerprinted the deceased and circulated copies to Australian states, New Zealand, and overseas fingerprint bureaus. A fingerprint chart reproduction exists as DOC-0016. The full bureau-response file and original search logs are not acquired.

Lawson made a cast of the features, and the bust was later marked as exhibit C.17. Embalming, formalin, and postmortem changes limit living-appearance inference. Biological specimens were submitted to Cowan in 1948, but retention or disposal status remains unknown.

\begingroup
\scriptsize
\setlength{\tabcolsep}{2pt}
\renewcommand{\arraystretch}{1.12}
\begin{longtable}{>{\raggedright\arraybackslash}p{0.12\textwidth}>{\raggedright\arraybackslash}p{0.18\textwidth}>{\raggedright\arraybackslash}p{0.34\textwidth}>{\raggedright\arraybackslash}p{0.24\textwidth}}
\caption{Identification material}\label{tab:chapter9-identification-material}\\
\toprule
ID & Type & Evidentiary value & Limitation \\
\midrule
\endfirsthead
\toprule
ID & Type & Evidentiary value & Limitation \\
\midrule
\endhead
FID-0001 & Dental chart & High for chart existence. & Exact notation needs specialist/manual transcription. \\
FID-0002 & Dental characteristics & High as identification feature. & No comparative dental records acquired. \\
FID-0003 & Fingerprints & High for fingerprint chart existence. & Original chart/search logs incomplete. \\
FID-0004 & Fingerprint circulation & High for process. & Full bureau-response file unavailable. \\
FID-0005 & Photographs & High for existence of full-face and side-face photographs. & Original exhibit images not fully catalogued. \\
FID-0006 & Death mask/bust & High for creation. & Postmortem changes limit living-appearance inference. \\
FID-0008 & Scars/marks & Moderate identifier value. & No complete body map acquired. \\
FID-0009 & Biological samples & High for 1948 submission. & Present retention/disposal unknown. \\
\bottomrule
\end{longtable}
\endgroup

\section{Modern Forensic Work}

The modern forensic layer is important but must be separated from the acquired 1948--1949 medical record. Mission 11 records the 2021 exhumation as an official South Australia Police and Forensic Science SA action linked to DNA recovery attempts. It also records researcher/lab-led DNA and genealogical work supporting the Carl Charles Webb attribution.

The repository does not contain a final coroner, South Australia Police, or Forensic Science SA identification report. It also does not contain a complete public DNA laboratory report or peer-reviewed DNA identification paper. The current identity wording therefore remains controlled by the identity-status guidance: researcher/lab-led Carl Charles Webb attribution, pending official coroner or South Australia Police determination.

Modern DNA, hair, and exhumation material do not resolve the medical cause of death in this chapter. They may affect identity assessment if official confirmation is acquired, but they do not change the 1948--1949 medical evidence unless linked to a documented modern forensic report.

\section{Medical Assessment}

The established observations are body condition, missing teeth, intense rigor, deep lividity, stomach congestion and blood mixed with food, organ congestion, normal but contracted heart appearance, spleen enlargement, specimen submission, negative common-poison testing, fingerprints, dental charting, photographs, and death-mask or bust creation.

Supported interpretations are narrower. The record supports that death was not explained by apparent natural disease in the reviewed testimony and that multiple medical witnesses considered poison or rare toxic agents. It does not support naming a detected poison, excluding all poisons, or closing cause of death.

The medical uncertainty register controls the unresolved issues: exact medical cause of death, whether any toxic agent was present, Hicks's exact drug group, reliability of the time-of-death estimate, place of death, stomach-content timing, possible missed injection marks, dental notation, biological-specimen fate, and liver-lobule significance.

\begin{table}[htbp]
\centering
\caption{Medical uncertainty register}
\label{tab:chapter9-medical-uncertainties}
\begin{tabularx}{\textwidth}{>{\raggedright\arraybackslash}p{0.14\textwidth}>{\raggedright\arraybackslash}p{0.42\textwidth}>{\raggedright\arraybackslash}p{0.20\textwidth}>{\raggedright\arraybackslash}X}
\toprule
ID & Question & Status & Priority \\
\midrule
MU-0001 & What was the exact medical cause of death? & Unresolved & Critical \\
MU-0002 & Was any toxic agent present? & Unresolved & Critical \\
MU-0003 & Which drug group did Hicks identify in exhibit form? & Unresolved & Critical \\
MU-0004 & How reliable was the 2 a.m. time-of-death estimate? & Partially documented & High \\
MU-0005 & Did the body die where found? & Unresolved & High \\
MU-0006 & Can stomach contents narrow the final meal window? & Partially documented & Medium \\
MU-0007 & Were injection marks reliably documented as absent? & Unresolved & High \\
MU-0008 & What are the exact dental identifiers? & Partially documented & Medium \\
MU-0009 & What was the present fate of biological specimens? & Unknown & High \\
MU-0010 & What was the medical significance of liver-lobule findings? & Unresolved & Medium \\
\bottomrule
\end{tabularx}
\end{table}
